% ============================================================
%  CAPÍTULO 7 — INNOVACIONES (Subdoc. 13)
%  Contenido: 5 innovaciones obligatorias + traza con arquitectura
% ============================================================
%  FUENTES:
%  - Art. 28° (Bases_Administrativas.md §469-481): cartera obligatoria de
%    cinco innovaciones, una por cada tipo obligatorio, sin repetir tipo.
%    Tipos: §475-481 (producto/servicio, proceso, tecnológica/arquitectura,
%    modelo de negocio/contratación, UX/sostenibilidad/impacto social).
%  - Art. 29° (§485-497): elementos exigidos por cada innovación (Formulario T-19):
%    problema concreto del caso que resuelve, tecnología que la sustenta, etc.
%    La omisión reduce la innovación a un enunciado.
%  - Subdoc. 13 (§1636-1641) + T-22 (§1842): presentación de la cartera.
%  - Pertinencia al caso y trazabilidad con arquitectura/EDT/flujo de caja (§28.2).
%  - APA 7.ª ed. exigida en tipo 3 (§479).
%  Ver FUENTES.md para el detalle completo.
\chapter{Innovaciones (Subdoc.\ 13)}

\nota{Las Bases Técnicas Transversales (Art. 3.3) exigen presentar 5 innovaciones, al menos 1 obligatoria por tipología: (1) producto o servicio, (2) proceso, (3) tecnológica o de arquitectura, (4) modelo de negocio o de contratación, (5) experiencia de usuario, sostenibilidad o impacto social. Cada innovación debe: estar alineada a las necesidades del caso, ser descrita en 2 a 4 párrafos, incluir un esquema o diagrama, y justificar por qué es innovadora (no solo una buena práctica). El evaluador busca creatividad técnica aplicada, no adjetivos vacíos.}

% -----------------------------------------------------------
\section{Innovación 1 — Producto o Servicio}
% -----------------------------------------------------------

\nota{Describir una innovación que el proponente incorpore a su propuesta de solución, que implique un cambio en el servicio entregado al cliente. Ejemplos del caso: un servicio de trazabilidad lot-a-lote con alertas automáticas para el ISP, un servicio de\"predicción de demanda por ruta\" que el cliente no pidió pero que resuelve un problema real. Debe ser concreta, medible y justificada.}

\subsection{Descripción}
\pendiente

\subsection{Por qué es innovadora}
\nota{Justificar: ¿qué tiene esta propuesta que no tiene la solución convencional del mercado? ¿Es un servicio nuevo? ¿Resuelve un problema que otros no ven? ¿Reduce costos de forma medible?}

\pendiente

\subsection{Esquema o Diagrama}
\pendiente

% -----------------------------------------------------------
\section{Innovación 2 — Proceso}
% -----------------------------------------------------------

\nota{Describir una innovación en un proceso del negocio del cliente. Ejemplos del caso: automatizar el retiro sanitario de un proceso que hoy toma días, reemplazar el conteo cíclico manual por un conteo continuo con dispositivos, digitalizar la rendición de efectivo. Debe ser un cambio real en la forma de operar, no solo una herramienta nueva sobre un proceso viejo.}

\subsection{Descripción}
\pendiente

\subsection{Por qué es innovadora}
\pendiente

\subsection{Esquema o Diagrama}
\pendiente

% -----------------------------------------------------------
\section{Innovación 3 — Tecnológica o de Arquitectura}
% -----------------------------------------------------------

\nota{Describir una innovación en la tecnología o la arquitectura del sistema. Ejemplos: uso de edge computing para operación desconectada del centro de distribución, IoT con sensores de temperatura de bajo costo, arquitectura event-driven para trazabilidad en tiempo real, uso de inteligencia artificial para optimización de rutas con restricciones reales del caso. Debe ser técnicamente justificada y viable con las tecnologías del Cap. 4.}

\subsection{Descripción}
\pendiente

\subsection{Por qué es innovadora}
\pendiente

\subsection{Esquema o Diagrama}
\pendiente

% -----------------------------------------------------------
\section{Innovación 4 — Modelo de Negocio o de Contratación}
% -----------------------------------------------------------

\nota{Describir una innovación en cómo se presta el servicio o cómo se contrata. Ejemplos: propuesta de tarifa variable por volumen procesado en vez de fija, modelo de shared services donde otros distribuidores de la zona pueden usar la plataforma (escalabilidad del negocio del proponente), propuesta de pago por resultado (porcentaje de reducción de mermas). Debe ser realista y viable económicamente.}

\subsection{Descripción}
\pendiente

\subsection{Por qué es innovadora}
\pendiente

\subsection{Esquema o Diagrama}
\pendiente

% -----------------------------------------------------------
\section{Innovación 5 — Experiencia de Usuario, Sostenibilidad o Impacto Social}
% -----------------------------------------------------------

\nota{Describir una innovación que mejore la experiencia de los usuarios finales (conductores, preventistas, bodegueros, clientes) o que tenga un impacto ambiental/social positivo. Ejemplos: app de conductores con gamificación para reducir mermas, reducción de huella de carbono por optimización de rutas, inclusión digital de preventistas que hoy usan cuadernos. Debe ser medible (encuesta de satisfacción, reducción de quejas, etc.).}

\subsection{Descripción}
\pendiente

\subsection{Por qué es innovadora}
\pendiente

\subsection{Esquema o Diagrama}
\pendiente

% -----------------------------------------------------------
\section{Traza de las Innovaciones con la Arquitectura}
% -----------------------------------------------------------

\nota{Tabla que cruce cada innovación con los componentes de la arquitectura (Cap. 4) que la habilitan, y los módulos del Cap. 3 que impacta. Esto demuestra que las innovaciones no son ideas sueltas sino que están integradas en la solución.}

\begin{tabularx}{\textwidth}{|L{2.5cm}|L{4cm}|L{4cm}|X|}
\hline
\rowcolor{lafroxClaro}
\textbf{Innovación} & \textbf{Tipología} & \textbf{Componente(s) de Arq.} & \textbf{Módulo(s) impactado(s)} \\
\hline
1 & Producto/Servicio & \pendiente & \pendiente \\
\hline
2 & Proceso & \pendiente & \pendiente \\
\hline
3 & Tecnológica & \pendiente & \pendiente \\
\hline
4 & Modelo de negocio & \pendiente & \pendiente \\
\hline
5 & UX/Sostenibilidad & \pendiente & \pendiente \\
\hline
\end{tabularx}
